% Abstract A of ABSTRACTS.md (descended from the review copy). B and C are alternatives.
\begin{abstract}
\noindent All 13.99 million USPTO trademark case files are re-parsed into
event-dated prosecution histories, and every filing is scored on how
surprising its goods/services language is against what its industry filed in
the five years before and the five years after. The average of the two
surprises is the filing's atypicality; their difference is its lead, positive
for wording the industry then moved toward. Among 3.10 million registrations,
the most leading fifth are cancelled at the five-year proof of continued
use 1.4 points more often than the most lagging fifth, within Nice class
and registration year with drafting and counsel held ($t = 8.3$). The penalty
holds under fifty global themes, five hundred, and fifty themes fitted per
class. It attaches to language new to the whole record: a theme established
elsewhere and extended into a new Nice class carries no penalty --- those
extensions survive the gate slightly better than the class's own
established themes, and they are most of what a global model scores as
atypical. The penalty was two to three
times larger for marks filed in the late 1990s, reversed for 2000--2004
filings, and returned from 2008, with the swing inside four technology
classes. Among one million owners debuting 2009--2018, the debut's language
does not predict a priced round; SEC reporting and listing fall in lead. The
firms holding most linked revenue described their first product in language
their industry already used. Applicants who refile an abandoned mark rewrite
toward the class's typical description, so every registered description is
in part a conformed one. Code, panels, and the corpus build are public at
\url{https://github.com/ericsilver/tm-vocabulary}.
\end{abstract}
